\documentclass[12pt, a4paper]{article}

\usepackage{amsmath}
\usepackage{amssymb}
\usepackage{amsthm}
\usepackage{geometry}
\geometry{margin=2.5cm}

\newtheorem{theorem}{Theorem}

\title{\textbf{Irrationality of $\sqrt{2}$}}
\author{0xHaru}
\date{\today}

\begin{document}

\maketitle

\begin{theorem}
    $\sqrt{2}$ is not rational.
\end{theorem}

\begin{proof}
Assume, for contradiction, that $\sqrt{2}$ is a rational number, meaning that
\[
    \exists\, a, b \in \mathbb{Z} : \frac{a}{b} = \sqrt{2}.
\]
If $a$ and $b$ have a common factor, it can be eliminated. Then $\sqrt{2}$ can be
written as an irreducible fraction $\dfrac{a}{b}$ such that $a$ and $b$ are
\emph{coprime}, i.e.\ they have no common factor other than $1$. This additionally
means that at least one of $a$ or $b$ must be odd.

It follows that
\[
    \frac{a^{2}}{b^{2}} = 2 \qquad \text{and therefore} \qquad a^{2} = 2b^{2}.
\]
Hence $a^{2}$ is even, because it equals $2b^{2}$, which is necessarily even as it
is two times another integer. Since squares of odd integers are never even, it
follows that $a$ itself must be even.

Because $a$ is even, there exists an integer $k$ such that $a = 2k$. Substituting
into $a^{2} = 2b^{2}$ gives
\[
    2b^{2} = a^{2} = (2k)^{2} = 4k^{2} \implies b^{2} = 2k^{2}.
\]
It follows that $b^{2}$ is also even, which in turn means that $b$ is even.

\medskip

So $a$ and $b$ are both even, but this contradicts the assumption that
$\dfrac{a}{b}$ is irreducible, since $a$ and $b$ would share the (non-trivial) common factor $2$.

\medskip

This means that $\sqrt{2}$ is not rational.
\end{proof}

\end{document}
